\section{Group attribution}
\label{sec:groups}

% DRAFT 1 (2026-10-20). Only the first subsection is written.

\todo{Opening paragraph: the idea of the fix. Give credit to groups of frames, chosen so that frames that show the same content fall in the same group, instead of to single frames.}

\subsection{Which attribution methods can be applied to groups}
\label{sec:which-groups}

Each of the methods in Section~\ref{sec:why} can assign a score to a group of frames $G \subseteq N$, but grouping means something different for each of them, and it removes dilution only for Shapley values.

\paragraph{Shapley values.}
The groups of a partition $\mathcal{G}$ become the players of a new game, with value $u(\mathcal{T}) = v\bigl(\bigcup_{G \in \mathcal{T}} G\bigr)$ for a set of groups $\mathcal{T} \subseteq \mathcal{G}$.
A group is then always kept or removed as a whole, and its Shapley value $\Psi_G$ is not divided among its frames.
This value is in general not the sum of the frame-level Shapley values of its frames: grouping changes how the credit is shared between the group and the other frames, not only within the group.
For example, with two identical frames $a_1, a_2$ and one frame $b$, and $v = 1$ only when $b$ and at least one of $a_1, a_2$ are present (and $0$ otherwise), the group $\{a_1, a_2\}$ receives $1/2$, while $\phi_{a_1} + \phi_{a_2} = 1/3$.
Play Fair computes the same quantity, so it can be grouped in the same way.
\todo{State and prove the second proposition here: if the model's output is the same whether the content is shown once or $m$ times, for every set of the other groups, then every group's Shapley value is unchanged; if the outputs differ by at most $\delta$, the values differ by at most $2\delta$. Proof in the appendix.}

\paragraph{Leave-one-out.}
Removing a whole group from the full clip gives the score $v(N) - v(N \setminus G)$, which is the importance $I(G)$ (Equation~\ref{eq:importance}).
If all frames that show a content lie in the same group, this score is not diluted.
But it is the quantity we use as the reference for the model's own dependence, so it cannot be evaluated against that reference.
It also does not decompose the output: the scores of the groups need not add up to $v(N) - v(\emptyset)$.
If frames that show the same content are split across two groups, removing one group leaves the other, and both scores are again close to zero.

\paragraph{Occlusion.}
Replacing every frame of a group with a flat frame is possible.
It differs from $I(G)$ only by the model's response to the inserted flat frames, which on our models is often as large as the importance itself (Appendix~\ref{app:flat-frame}).

\paragraph{Integrated Gradients and Grad-CAM.}
For these methods, the natural score of a group is the sum of the scores of its frames.
For Integrated Gradients this is exact, because the scores of all input elements add up by completeness.
So grouping adds nothing to the frame-level scores: the score of a group of copies is the sum of the scores of the copies.
This sum grows with the number of copies whenever each copy keeps its score, so it does not follow the model's dependence on the content either.

When each frame displays the score of its group, the score of a part of the video is the score of its group, not the sum of the scores displayed on its frames: adding up the displayed scores would multiply the group's score by its number of frames.

\todo{Remaining subsections: how the groups are found (DINOv2 similarity, average-linkage threshold), and how group scores are shown per frame (broadcast scores).}
